\clearpage
\VerdierPhysicalPageAnchor{0241}%
\begingroup\small\setlength{\parskip}{0pt}\setlength{\emergencystretch}{1.5em}
\setlength{\abovedisplayskip}{4pt plus 2pt minus 2pt}%
\setlength{\belowdisplayskip}{4pt plus 2pt minus 2pt}%

\VerdierSemanticAnchor{ch4d-4-5-2}
\noindent\textbf{4.5.2.}\quad
$F\colon D(\mathcal A)\longrightarrow\mathcal B$ を右定常な
コホモロジー的関手、$p$ を整数、$X$ を
$\operatorname{comp}(\mathcal A)$ の対象とし、

\VerdierSemanticAnchor{ch4d-4-5-2-second-filtration}
\[
0=X(-\infty)\longrightarrow\cdots\longrightarrow X(q)
\longrightarrow X(q+1)\longrightarrow\cdots
\longrightarrow X(+\infty)=X
\]
を $X$ の第二標準フィルトレーション (4.3.1) とする。帰納系
$F^p(X(q))$ は定常である。すなわち、ある整数 $q_0$ が存在し、
$q_0\leq q\leq+\infty$ を満たすすべての $q$ に対して射
\[
F^p(X(q))\longrightarrow F^p(X)
\]
は同型である。実際、(4.2.3) の記法を用いると完全列
\[
F^{p-1}\bigl(X(q,+\infty)\bigr)\longrightarrow F^p(X(q))
\longrightarrow F^p(X)\longrightarrow F^p\bigl(X(q,+\infty)\bigr),
\]
がある。関手 $F$ は右定常だから、ある整数 $q_0$ が存在し、すべての
$q>q_0$ に対して対象 $F^p(X(q,+\infty))$ と
$F^{p-1}(X(q,+\infty))$ は零である。

\VerdierSemanticAnchor{ch4d-4-5-2-first-filtration}
同様に、$F\colon D(\mathcal A)\longrightarrow\mathcal B$ を右定常な
コホモロジー的関手、$p$ を整数、$Y$ を
$\operatorname{comp}(\mathcal A)$ の対象とし、
\[
0=Y(-\infty)\longrightarrow\cdots\longrightarrow Y(q)
\longrightarrow Y(q+1)\longrightarrow\cdots
\longrightarrow Y(+\infty)=Y,
\]
を $Y$ の第一フィルトレーション (4.1.2) とする。ある整数 $q_0$ が存在し、
すべての $q<q_0$ に対して対象 $F^p(Y(q))$ は零である。先と同様の検証は
読者に委ねる。読者は同じく反対圏へ移ることにより、左定常な関手の類似の
性質を調べることができる。

\VerdierSemanticAnchor{ch4d-4-5-3}
\medskip\noindent\textbf{命題 4.5.3.}\quad

\VerdierSemanticAnchor{ch4d-4-5-3-a}
\noindent\textit{a) $Y$ を $\operatorname{comp}^{*}(\mathcal A)$ の対象
(1.1.1)（$*=+,-,b,{}$ ``空''）とする。
$F\colon D^*(\mathcal A)\longrightarrow\mathcal B$ をコホモロジー的関手
とする。複体 $Y$ に関する関手 $F$ の第一スペクトル系列 (4.4.1) の初期項は
(4.4.2)：}
\[
\textit{$\mathrm{I}_1^{p,q}=F^q(Y^p)$}
\]
\textit{であり、次の場合に定常である（第 II 章、4.4.2）：}

\begin{enumerate}
\item[1)] \textit{$Y$ は $\operatorname{comp}^{b}(\mathcal A)$ の対象である。}
\item[2)] \textit{$Y$ は $\operatorname{comp}^{+}(\mathcal A)$ の対象であり、
$F$ は右定常である。}
\item[3)] \textit{$Y$ は $\operatorname{comp}^{-}(\mathcal A)$ の対象であり、
$F$ は左定常である。}
\item[4)] \textit{関手 $F$ は定常である。}
\end{enumerate}

\VerdierSemanticAnchor{ch4d-4-5-3-b}
\noindent\textit{b) $X$ を $D^*(\mathcal A)$ の対象 (1.2.4)、
$F\colon D^*(\mathcal A)\longrightarrow\mathcal B$ をコホモロジー的関手
とする。複体 $X$ に関する関手 $F$ の第二スペクトル系列 (4.4.5) は、
次の場合に定常である：}
\endgroup
